\documentclass[10pt]{article}
\usepackage[margin=20mm]{geometry}
\usepackage{amsmath,amssymb,fontspec,unicode-math,xurl,hyperref}
\IfFontExistsTF{DejaVu Serif}{\setmainfont{DejaVu Serif}}{\setmainfont{Latin Modern Roman}}
\setmathfont{latinmodern-math.otf}
\setlength{\parindent}{0pt}\setlength{\parskip}{6pt}\sloppy
\begin{document}
Abstract: The printed Gaussian variational model, ELBO and posterior-predictive identities are valid in the inspected scope. The released implementation has two core defects: local inference accumulates shared-decoder gradients, including validation gradients, into a subsequent global update; and final perplexity reuses the random noise to which the variational parameters were fitted. A Gaussian-prior, smooth categorical example shows that the latter estimation rule can violate the claimed perplexity-bound interpretation even in expectation. Both are assessed Major for the pinned implementation. The release is later than the paper, so their numerical effect on historical results remains unverified. A per-token compute formula also retains a per-sequence factor. No author program was executed and no historical score failure is invented.\par
Author: AI-assisted technical review\par
Date:\par
Title: Technical Verification of Latent Thought Model Inference and Evaluation\par
\# Identity and source versions\par
Deqian Kong, Minglu Zhao, Dehong Xu, Bo Pang, Shu Wang, Edouardo Honig, Zhangzhang Si, Chuan Li, Jianwen Xie, Sirui Xie and Ying Nian Wu, Latent Thought Models with Variational Bayes Inference-Time Computation, ICML 2025, PMLR 267 [Retained latex: \textbackslash{}cite\{paper\}]. Assignment 0218 / I25-43587; OpenReview xCWAcX4pNa. The official publisher PDF retrieved 25 October 2026 has 18 pages, with appendices on pp. 13–18. All source page numbers below are one-based PDF pages and agree with printed local numbers.\par
The official project page identifies the released repository . The prepared snapshot is commit , dated 30 October 2025 [Retained latex: \textbackslash{}cite\{code\}]. Its nine files were preserved completely, without omitted members, and their staged hashes and archive/Git-blob mappings checked. This commit postdates the conference paper. No historical run manifest or original experiment revision binds it to the published numbers.\par
The parent inspected the original model and evaluation descriptions, the full posterior optimizer, the relevant decoder/ELBO and training paths, and configuration; independently performed the calculations below; and viewed source pages 4, 5 and 7. Three native auxiliaries independently checked mathematics, source control flow and counterreadings. No author code, dataset, training job, judge or paid API was executed. The snapshot has no separate zero-shot, GSM8K or MAUVE evaluation driver or historical result artifacts. Optional OpenReview access returned HTTP 403, and later revisions/errata were not retrieved.\par
\# Definitions and mathematical checks that survive\par
The model is \[p_\beta(x)=\int p(z)p_\beta(x\mid z)\,dz,\qquad
 p(z)=\mathcal N(0,I),\qquad
 p_\beta(x\mid z)=\prod_{n=1}^N
 p_\beta(x^{(n)}\mid x^{(<n)},z).\] Each sequence has local diagonal-Gaussian variational parameters \(q_\psi(z\mid x)=\mathcal N(\mu,\operatorname{diag}(\sigma^2))\). The decoder parameters \(\beta\) are shared. Equation (5), p. 4, is \[\mathcal L(\beta,\psi;x)
 =\mathbb E_{q_\psi}\log p_\beta(x\mid z)
  -D_{\rm KL}(q_\psi\|p).\] The identity \[\log p_\beta(x)-\mathcal L(\beta,\psi;x)
 =D_{\rm KL}(q_\psi\|p_\beta(z\mid x))\ge0\] proves the ELBO bound whenever the expectations are well-defined. Optimizing \(\psi\) using the entire observed test sequence is legitimate for density estimation under this identity; it is not itself evidence of answer leakage.\par
The Fisher identity in Eq. (3), the Langevin noise scaling in Eq. (4), and the conditional predictive integral in Eq. (6) are also consistent under their usual differentiation/integration assumptions. Conditional generation infers \(q(z\mid x)\) from the prefix \(x\), then generates the completion; the formula does not require access to the true completion. Finite-step inference bias is explicitly acknowledged on p. 4.\par
Importantly, the released charges the full Gaussian KL, \[K=\frac12\sum_j(\mu_j^2+\sigma_j^2-1-\log\sigma_j^2),\] with coefficient one. It sums the reconstruction loss over tokens, adds \(K\), divides by sequence length, exponentiates per sequence, and averages the resulting perplexities across the batch. The issue is not omission or arbitrary downweighting of KL. The validation caller records this perplexity separately from a reconstruction-only loss.\par
\# Core implementation findings\par
Problem 1 (Local inference and validation gradients enter the shared-decoder update).  The pinned training path does not isolate fast local optimization from slow global optimization.\par
\#\# Source claim and complete control-flow evidence\par
Algorithm 1, p. 4, separates updating each sequence’s local \((\mu,\sigma^2)\) from a slow update of the shared decoder. During validation, local fitting should leave shared parameters and their training optimizer state unaffected.\par
In , the inner AdamW optimizer owns only `mu` and `log\_var`. Its `zero\_grad` at line 117 therefore clears only those variables. The backward calls at lines 122/126 traverse the decoder as well as the local variables. Decoder weights retain their gradient requirements; changing the module to evaluation mode does not freeze them. The code detaches an optional cache, not the full decoder computation.\par
The validation routine in repeatedly calls this same posterior optimizer on `val` batches. Periodic evaluation at lines 336–338 precedes the training micro-batches. Their inner backward passes and final reconstruction backward pass add more gradients. The global optimizer consumes the accumulated decoder gradients at line 398. Only line 402 clears them, after the step. The inspected source contains no intervening decoder freeze or reset that cancels this path.\par
For the ordinary single-GPU branch, this yields \[g_{\rm used}
 =g_{\rm intended\ slow}
  +\sum_{\rm training\ fast\ steps}g_{\beta,{\rm fast}}
  +\sum_{\rm validation\ fast\ steps}g_{\beta,{\rm val}},\] before clipping and the optimizer transformation. The last sum is present on evaluation iterations. The sums also have different normalization from the intended slow loss: inner ELBO losses are summed over tokens/batches, whereas the final slow cross-entropy is averaged and divided for gradient accumulation.\par
The actual default configuration has 16 fast steps, 100 validation batches, evaluation every 1,000 iterations, and a positive global learning rate at iteration 1,000. It uses `bfloat16`, so the gradient scaler is disabled. This ordinary supported path is sufficient; no behavior of a particular multi-GPU execution is assumed.\par
\#\# Independent calculus check and counterreadings\par
The mechanism does not depend on a large Transformer. For a scalar decoder logit \(\beta+\mu\) and a validation target one, the negative log likelihood is \(-\log\sigma(\beta+\mu)\). At \(\beta=\mu=0\), \[\partial_\beta\ell_{\rm val}=-\tfrac12.\] An optimizer that clears only \(\mu\) leaves that decoder gradient untouched. Even a subsequent training loss with zero decoder gradient cannot remove it by addition.\par
Zero global learning rate on the first iteration is not a general rescue. With the supplied Adam coefficients \((0.9,0.95)\), one such gradient creates first/second moments \(-0.05\) and \(0.0125\) even if that step’s learning rate is zero. At the next step, with zero new gradient and learning rate \(4\times10^{-7}\), the bias-corrected Adam update moves an initially zero parameter by approximately \(2.71459\times10^{-7}\). Without the validation gradient it remains zero. This is an independent analytic optimizer calculation, not a run of the author software. Later evaluation steps already have positive learning rates and establish the problem without relying on this moment-state observation.\par
An evaluation-only script that never makes another global update would not have this training contamination. The finding concerns the released combined training/validation loop. Nor is it an accusation that legitimate local posterior fitting inherently trains the global model on test answers; it is the specific uncancelled global-gradient path that does so here.\par
\#\# Impact and precise repair boundary\par
This is Major for the released core learning/evaluation implementation: the slow update differs from the stated dual-rate algorithm, and held-out validation targets can affect shared training. It directly threatens the interpretation of validation perplexity and scaling curves generated by this path. The exact magnitude and whether the historical published runs used this version are Needs Evidence; no published score is declared refuted on provenance alone.\par
A correct isolation rule is to compute fast-step gradients only for the local variational parameters, or temporarily freeze decoder parameters while retaining differentiation through the decoder to its latent inputs. Restore the decoder’s gradient flags for the slow step. Validation must leave both decoder gradients and global optimizer state unchanged. Deliberate outer micro-batch gradients must be preserved, so indiscriminately clearing all gradients after every inner step is not a general repair. These rules remove the identified paths; implementation testing and retraining would still be needed to validate revised results. No author source was modified.\par
Problem 2 (Fitting and scoring with the same noise invalidates the fitted-ELBO estimator).  The released optimizer uses one random tensor for posterior fitting and final perplexity, unlike Algorithm 1’s sampling inside the fast loop.\par
\#\# Actual source and the statistical distinction\par
draws \(e\sim\mathcal N(0,I)\) once, before the loop at line 112. Every ELBO call at line 119 uses that same \(e\). The final latent at line 141 and final metric at line 145 reuse it again. Thus \(\widehat\psi\) depends on \(e\) and the scored point is \[z(e)=\mu_{\widehat\psi(e)}
      +\sigma_{\widehat\psi(e)}\odot e.\] It is not a fresh draw from the fitted Gaussian conditional on \(\widehat\psi\).\par
Common random numbers can be a legitimate optimization device. Also, even an independent one-sample ELBO estimate need not be a pointwise bound on a particular draw. The decisive issue here is stronger: adaptive reuse can invalidate the bound interpretation even after averaging over the scoring randomness. Table 1, p. 7, explicitly labels LTM perplexities as upper bounds, and the released validation metric is offered as such an ELBO-derived statistic.\par
\#\# A fully specified Gaussian/categorical counterexample\par
Consider the probabilistic family in Eqs. (1), (2) and (5), specialized to one binary token: \[z\sim\mathcal N(0,1),\qquad
 p(x=1\mid z)=\sigma(kz),\qquad k=100.\] This is a smooth two-class softmax decoder. Symmetry and \(\sigma(kz)+\sigma(-kz)=1\) give \(p(x=1)=1/2\), hence true perplexity two. Take the observed token to be one and draw the fixed reparameterization noise \(e\sim\mathcal N(0,1)\).\par
Let \(a=0.1\) and choose the following admissible fitted Gaussian for each noise draw: \[(\mu(e),s(e))=
 \begin{cases}
 (a-e/\sqrt{1+e^2},\,1/\sqrt{1+e^2}),&e<a,\\
 (0,1),&e\ge a,
 \end{cases}
 \qquad q_e=\mathcal N(\mu(e),s(e)^2).\] Every variance is strictly positive and every KL is finite. The reparameterized scoring point is \(z=\mu(e)+s(e)e=\max(e,a)\), so \[-\log p(x=1\mid z)\le
 \epsilon:=\log(1+e^{-ka}).\] For \(e<a\), direct substitution into the full Gaussian KL gives \[K(e)=\tfrac12\left[
 a^2-\frac{2ae}{\sqrt{1+e^2}}+\log(1+e^2)
 \right]
 \le a+\tfrac12a^2+\tfrac12\log(1+e^2).\] For \(e\ge a\), \(K(e)=0\). No KL term is omitted.\par
Write \(\phi,\Phi\) for the standard normal density and distribution. The mean of the reported one-sample perplexity satisfies \[\begin{aligned}
 \mathbb E\widehat{\rm PPL}
 &\le (1+e^{-ka})
 \left[1-\Phi(a)+
 e^{a+a^2/2}\mathbb E\{\sqrt{1+e^2}\mathbf1_{e<a}\}\right]\\
 &\le (1+e^{-ka})
 \left[1-\Phi(a)+e^{a+a^2/2}
 \sqrt{\Phi(a)\{2\Phi(a)-a\phi(a)\}}\right]\\
 &=1.292448567\ldots<2.\end{aligned}\] The second line uses Cauchy–Schwarz and \(\mathbb E[e^2\mathbf1_{e<a}]=\Phi(a)-a\phi(a)\). Thus the estimator can report a value below true perplexity even in expectation.\par
The construction extends to any sequence length, including \(N=1024\): use independent prior coordinates, independent binary conditional factors, the observed all-one sequence, and the product of the above Gaussian variational factors. The true marginal is \(2^{-N}\) and true per-token perplexity remains two. The released scoring functional exponentiates summed sampled NLL plus KL divided by \(N\), which is the geometric mean of the one-coordinate perplexities. By the arithmetic–geometric mean inequality, its expectation is still at most \(1.292448567\). This also addresses the code’s actual order of operations: it averages per-sequence perplexities rather than exponentiating a batch-averaged loss.\par
All probabilistic assumptions behind the ELBO argument are satisfied: a normalized standard Gaussian prior, a normalized smooth categorical likelihood, positive diagonal Gaussian variances, and the complete KL. The example is a counterexample to the general estimation rule, not a simulation of the trained Transformer or a claim that its finite AdamW trajectory returns this particular family. An exact optimizer of the sampled objective could do at least as well on that objective as this feasible family, so the defect is not cured merely by longer optimization.\par
\#\# Independent finite-step numerical corroboration\par
A separate parent-written scalar calculation uses \(p(x=1\mid z)=\sigma(8z)\) and 16 AdamW steps, with \(\mu_0=0\), \(\log s_0^2=-5+0.1\eta\), independent standard-normal \(\eta,e\), coefficients \((0.9,0.999)\), weight decay \(0.01\), denominator offset \(10^{-8}\), and step sizes \[\gamma_t=0.034+\frac{0.34-0.034}{2}
 [1+\cos(\pi t/16)],\qquad t=0,\ldots,15.\] Only the local Gaussian parameters are updated in this diagnostic. Deterministic Gaussian quadrature of orders 80 and 160 gives reused-noise mean negative ELBOs \(0.6693841\) and \(0.6692426\), respectively, below the true \(-\log p(x)=\log2=0.6931472\). Fresh-noise evaluation of the same fitted distributions gives \(1.0228832\) and \(1.0226480\). An independent auxiliary Monte Carlo calculation is consistent with these values.\par
These are numerical integrations of an independently written analytic-gradient toy, not rigorous quadrature error bounds or author-code execution. They corroborate the finite-step mechanism but do not reproduce the Transformer or historical scores. The preceding analytic construction remains the decisive general counterexample.\par
\#\# Impact, repair and preserved counterevidence\par
This is Major for the released core estimator: the implementation does not justify interpreting its adaptively reused-noise score as the ELBO-based perplexity upper bound used to support sample/parameter efficiency and inference-step scaling. The exact mathematical ELBO in Eq. (5) remains valid. Neither the example nor static inspection establishes the actual corrected value of the published \(3.05\), Table 1 entries, or every scaling curve.\par
Use independent scoring noise after fitting \(\psi\), or sufficiently accurate integration of the fitted expectation. Conditional on the fitted parameters, fresh Monte Carlo log-likelihood samples then estimate the actual ELBO without this adaptive reuse bias. Resampling within the fast loop also aligns the implementation with Algorithm 1. A finite Monte Carlo estimate is not a deterministic pointwise bound; precision and uncertainty still need reporting. No corrected experimental result is claimed.\par
More inference can also tighten a valid variational bound for a fixed decoder without changing its true marginal distribution. That ordinary distinction should be retained when discussing inference-step scaling; it is separate from the adaptive-noise defect above.\par
Problem 3 (The expression labeled per-token compute retains a sequence-length factor).  Section 2.4, p. 5, changes the stated unit without dividing the preceding per-sequence cost by \(N\).\par
The displayed forward cost for a length-\(N\) sequence is \[O\!\left(L(N^2H+NN_zH+NH^2)+NVH\right).\] Multiplying by the number of fast/slow passes preserves a per-sequence unit. Under that same dense-attention accounting, the per-token expression is \[O\!\left((T_{\rm fast}+1)
 \{L(NH+N_zH+H^2)+VH\}\right),\] obtained by division by \(N\). Here \(N_z\) follows Section 2.4’s explicitly stated per-layer convention, unlike the total-latent convention elsewhere. Efficient fixed-window attention would further change the attention term, but is not needed to identify this unit discrepancy.\par
This is Minor as a local cost expression. At a fixed training length, an omitted common factor does not by itself reverse configuration rankings, and no evidence binds the actual Table 1 numerical FLOP counts to evaluation of that erroneous symbolic expression. The report does not claim the published FLOP numbers are all inflated by \(1024\) or that the empirical compute comparisons have been numerically refuted.\par
\# Empirical claims: inspected scope and limits\par
The paper clearly separates the 3B-token pretraining example from Table 1’s equal-total-compute checkpoints, including the 0.9B-token LTM-Large checkpoint. Their different training-token counts are not inherently contradictory. The release configuration, however, uses width 768 and 12 heads, while Appendix A.1 specifies width 512 and eight heads for the reported models. This reinforces the need for a historical configuration binding rather than silently treating the current defaults as the published runs.\par
Appendix A.3, p. 13, discloses GSM8K pass@5 on a 1K test subset with at most 50 new tokens per candidate. GPT-2 uses beam width five, while LTM uses five posterior fits and multinomial samples. That is a decoding/protocol difference limiting attribution to architecture alone. No released caller establishes access to true test answers, and this review finds no demonstrated GSM8K answer leakage.\par
Conditional generation uses 1,000 50-token prefixes, five 50-token completions each, and borrowed baseline MAUVE values. The absence of seed-level outputs and confidence intervals limits claims about small score differences. The reported unconditional evaluation uses 64 length-1024 sequences, greedy decoding and GPT-2-XL scoring. Prior sampling for unconditional generation legitimately avoids posterior inference; it would be incorrect to charge all \(T_{\rm fast}\) steps to that task’s generation speed solely by analogy with conditional inference.\par
Progressive latent-layer inclusion on 200 validation sequences is evidence of reconstruction patterns (pp. 17–18). A causal interpretation as unique semantic synthesis layers would require controls separating layer position from the amount of latent information supplied, such as matched-count or alternative-order interventions. The present review does not disprove the reported reconstruction pattern. The VAE/Langevin comparisons likewise do not receive a new failure judgment without raw diagnostics and evaluation drivers.\par
None of these unresolved empirical questions is used to inflate the severity of the two demonstrated implementation findings. The original finite-step-bias discussion, correctly normalized KL, valid ideal probabilistic identities and explicit evaluation differences remain counterevidence against broader accusations.\par
\# Conclusion and verification status\par
Object | Judgment | Verified correction or remaining work\par
Decoder-gradient isolation | Major | Static training/validation path proved; isolate local gradients and retrain to assess scores.\par
Adaptive-noise perplexity | Major | Estimator counterexample; fresh fitted-posterior scoring restores the expectation identity.\par
Per-token compute expression | Minor | Divide sequence expression by \(N\); actual FLOP tables not refuted.\par
Ideal ELBO/predictive model | Scoped NIF | Identities and full-KL normalization checked.\par
Historical empirical values | Needs Evidence | Original code/config/checkpoints and independent evaluation needed.\par
The overall supported assessment is Major for the pinned released implementation. The scientific consequences for historical paper scores remain conditional on provenance. Repair verification here is limited to the stated mathematical/control-flow mechanisms; no corrected implementation or experiment was run. Adjudication remains pending, human reviews are empty, and author adoption has not been requested. No claim of first discovery or external erratum status is made.\par
\# Independent calculations and provenance\par
Original PDF SHA-256: . Source manifest SHA-256: . Preserved source ZIP SHA-256: . Key member SHA-256 values are:\par
- : .\par
- : .\par
- : .\par
All checked original/archive/member bytes remain unchanged. The longest checked source path is 112 characters; detailed mappings and auxiliary records are preserved.\par
Parent-authored independently evaluates the Gaussian moment bounds, the 1.292448567 perplexity bound and the Adam moment-state example. It also verifies a simpler fixed-variance diagnostic: \(\mu(e)=(0.1-e)_+\) gives expected sampled ELBO at least \(-0.292506084\), above the true \(-\log2\). Its 1024-coordinate moment-generating-function extension gives expected perplexity at most \(1.340027831\). These are analytic examples and ordinary numerical evaluation of closed forms, not author-code or Transformer execution. Actual outputs are in ; the independent finite-step calculation and quadrature settings are preserved in and .\par
The frozen mathematical-journal template is used, with 11-point serif text and conventional equations. Parent owns actual two-pass `pdflatex` compilation with shell escape disabled and visual inspection of every rendered final page. Commands, exit codes, exact final source/PDF hashes, logs and inspection evidence are in , , , , and .\par
Task-selected model label is 6 Pro; exposed parent system label is GPT-6 Astra Pro. Auxiliary model identifiers, visible UI label, conversation URL and accessible Space references were not exposed and are not invented. This is an initial review, without acceptance, publication or author contact.\par
9 D. Kong et al., Latent Thought Models with Variational Bayes Inference-Time Computation, PMLR 267 (2025). Official publisher PDF, frozen retrieval 25 October 2026. https://proceedings.mlr.press/v267/kong25c.html. M. Zhao et al., Latent-Thought-Language-Models, commit , 30 October 2025. https://github.com/mingluzhao/Latent-Thought-Language-Models.\par
\end{document}
